\documentclass[xelatex,mathserif,serif,notheorems]{beamer} % ,notheorems
\mode<presentation>{
  %% BEAMER CONFIG %%
  \usetheme{Madrid}
  \setbeamertemplate{navigation symbols}{}
  \setbeamersize{text margin left=0.5cm,text margin right=0.5cm}
  \usefonttheme[onlymath]{serif}
}

% \usepackage[latin1]{inputenc}
\usepackage[english]{babel}
% \usepackage{a4}
\usepackage{latexsym}
\usepackage{amssymb}
\usepackage{amsmath}
\usepackage{epsfig}
\usepackage[T1]{fontenc}
% \usepackage{mathptmx} % Use Times New Romans as font?
\usepackage{color}
\usepackage{epstopdf}
\usepackage{microtype}
\usepackage{hyperref}
\usepackage[useregional]{datetime2}
\DTMlangsetup[en-US]{showdayofmonth=false}
\usepackage{lipsum}

\usepackage{marvosym}

\renewcommand*\sfdefault{lmss}
\renewcommand*\ttdefault{txtt}

%Information to be included in the title page:
\title{(Q)\texttt{M}-types and Coinduction in HoTT / CTT}
\subtitle{Master's Thesis, Computer Science}
\author{\large Lasse Letager Hansen, 201508114\\[1ex]{\small Supervisor: Bas Spitters}}
\institute{Aarhus University}
\date{\today}

\renewcommand{\insertshortauthor}{Lasse Letager Hansen}
\titlegraphic{
  \epsfig{file=logo.eps}
}

% Custom packages added:

% \usepackage{fullpage}
% \setlength\marginparwidth{0.7in}

\usepackage{todonotes} % \todo command

\usepackage{proof} % \infer command

\usepackage{tikz-cd}
\usepackage{xfrac} % \sfrac - sideways fraction


% \usepackage{mathtools} % multlined

% allow breaks in align and equation (a new page for example)
\allowdisplaybreaks

\renewcommand{\thefootnote}{\fnsymbol{footnote}} %footnote style

\usepackage{amsthm}

% \newtheoremstyle{plain}
%   {\abovedisplayskip}   % ABOVESPACE
%   {\belowdisplayskip}   % BELOWSPACE
%   {\itshape}  % BODYFONT
%   {0pt}       % INDENT (empty value is the same as 0pt)
%   {\bfseries} % HEADFONT
%   {.}         % HEADPUNCT
%   {5pt plus 1pt minus 1pt} % HEADSPACE
%   {}          % CUSTOM-HEAD-SPEC
\theoremstyle{plain} % default
\newtheorem{thm}{Theorem}[section]
\newtheorem{lem}[thm]{Lemma}
\newtheorem{prop}[thm]{Proposition}
\newtheorem{cor}{Corollary}

\theoremstyle{definition}
\newtheorem{defn}[thm]{Definition}% [section]
\newtheorem{conj}{Conjecture}[section]
\newtheorem{exmp}{Example} % [section]

\theoremstyle{remark}
\newtheorem*{rem}{Remark}
% \newtheorem*{note}{Note} % Not used?
\newtheorem{case}{Case}

% Custom commands

\newcommand*{\thmref}[1]{\hyperref[thm:#1]{Theorem~\ref*{thm:#1}}} % Theorem~\ref
\newcommand*{\lemref}[1]{\hyperref[lem:#1]{Lemma~\ref*{lem:#1}}}
\newcommand*{\corref}[1]{\hyperref[cor:#1]{Corollary~\ref*{cor:#1}}}
\newcommand*{\defref}[1]{\hyperref[defn:#1]{Definition~\ref*{defn:#1}}}
\newcommand*{\figref}[1]{\hyperref[fig:#1]{Figure~\ref*{fig:#1}}}
\newcommand*{\tableref}[1]{\hyperref[table:#1]{Table~\ref*{table:#1}}}
\newcommand*{\exmpref}[1]{\hyperref[exmp:#1]{Example~\ref*{exmp:#1}}}
\newcommand*{\sectionref}[1]{\hyperref[sec:#1]{Section~\ref*{sec:#1}}}
\newcommand*{\subsectionref}[1]{\hyperref[sec:#1]{Subsection~\ref*{sec:#1}}}
\newcommand*{\chapterref}[1]{\hyperref[ch:#1]{Chapter~\ref*{ch:#1}}}

\newcommand*{\set}[1]{\left\{#1\right\}}
\newcommand*{\term}[1]{\textcolor{green!30!black}{#1}} % \textcolor{green!70!black} \textcolor{blue!80!black}
\newcommand*{\pathterm}[1]{\textcolor{green!65!red!75!black}{#1}}
\newcommand*{\type}[1]{\textcolor{magenta!90!black}{#1}}
% \newcommand*{\dependenttype}[1]{\textcolor{magenta!70!black}{#1}}
\newcommand*{\container}[1]{\textcolor{orange}{#1}}
\newcommand*{\containerpair}[2]{\ensuremath{\colorlet{savedleftcolor}{.}\color{orange}\left(\color{savedleftcolor}#1\,\textcolor{orange}{\mathbf{,}}\,#2\color{orange}\right)\color{savedleftcolor}}}
\newcommand*{\containerpairsimple}[2]{\containerpair{#1}{\lambda\,\_,\,#2}}
\newcommand*{\universe}[1]{\textcolor{orange!80!black}{#1}}
\newcommand*{\unit}{\type{\textbf{1}}}
\newcommand*{\empt}{\type{\textbf{0}}}
\newcommand*{\coalg}[2]{#1\texttt{-}#2}

\newcommand*{\relation}[1]{\textcolor{gray!75!black}{\ensuremath{\mathtt{#1}}}}
\newcommand*{\constant}[1]{\textcolor{orange!60!black}{\ensuremath{\mathtt{#1}}}}

\newcommand*{\function}[1]{\textcolor{blue!60!black}{\ensuremath{\mathtt{#1}}}}
\newcommand*{\constructor}[1]{\textcolor{purple!60!black}{\ensuremath{\mathtt{#1}}}}
\newcommand*{\destructor}[1]{\textcolor{yellow!60!black}{\ensuremath{\mathtt{#1}}}}
\newcommand*{\typeformer}[1]{\ensuremath{\mathtt{#1}}}
\newcommand*{\functor}[1]{\ensuremath{\mathbf{\mathtt{#1}}}}

\newcommand*{\unitelem}{\constant{\star}} % \constant{tt}
\newcommand*{\quotientconstructor}[1]{\constructor{[}\,#1\,\constructor{]}}

\newcommand*{\ssfrac}[2]{ {\LARGE \text{\sfrac{\ensuremath{#1}}{\ensuremath{#2}}}}}

\newcommand*{\natcases}[2]{\ensuremath{\mathbb{[\hspace{-.5mm}\rangle}~\ensuremath{#1}~,~\ensuremath{#2}~\mathbb{\langle\hspace{-.5mm}]}}}

% Concatenation of paths taken from HoTT book
\newcommand{\ct}{%
  \mathchoice{\mathbin{\raisebox{0.5ex}{$\displaystyle\centerdot$}}}%
             {\mathbin{\raisebox{0.5ex}{$\centerdot$}}}%
             {\mathbin{\raisebox{0.25ex}{$\scriptstyle\,\centerdot\,$}}}%
             {\mathbin{\raisebox{0.1ex}{$\scriptscriptstyle\,\centerdot\,$}}}
}

\newcommand*{\sym}[1]{\ensuremath{#1^{-1}}}

\newcommand{\setlengths}{
  \setlength{\abovedisplayskip}{4pt}
  \setlength{\belowdisplayskip}{4pt}
  \setlength{\abovedisplayshortskip}{2pt}
  \setlength{\belowdisplayshortskip}{2pt}
}
\newcommand{\startwitheq}{\vspace{-2.5mm}} % -\baselineskip

\begin{document}

\frame{\titlepage}

\section{Introduction}

\begin{frame}
  \frametitle{Overview}
  \tableofcontents[subsubsectionstyle=hide]
\end{frame}

\subsection{Background Theory}
\subsubsection{Inductive / Coinductive definitions}
\begin{frame}
  \frametitle{Inductive / Coinductive definitions}
  \framesubtitle{Inductive definition}
  We can define the natural numbers \emph{inductively} from the two constructors
  \begin{defn}[The Natural Numbers]\setlengths\startwitheq
    \strut
    \hfill
    \begin{minipage}{0.15\linewidth}
      \begin{equation}
        \vcenter{\infer{0 : \mathbb{N}}{\strut}}
      \end{equation}
    \end{minipage}
    \hfill
    \begin{minipage}{0.225\linewidth}
      \begin{equation}
        \vcenter{\infer{\constructor{succ}~n : \mathbb{N}}{n : \mathbb{N}}}
      \end{equation}
    \end{minipage}
    \hfill
    \strut
  \end{defn}
  for which we define an \emph{inductive} equivalence relation
  \begin{defn}[Equivalence for the Natural Number]\setlengths\startwitheq
    \strut
    \hfill
    \begin{minipage}{0.225\linewidth}
      \begin{equation}
        \vcenter{\infer[\constructor{\sim_{0}}]{0 = 0}{\strut}}
      \end{equation}
    \end{minipage}
    \hfill
    \begin{minipage}{0.45\linewidth}
      \begin{equation}
        \vcenter{\infer[\constructor{\sim_{succ}}]{\constructor{succ}~n~\relation{\sim_{\mathbb{N}}}~\constructor{succ}~m}{n~\relation{\sim_{\mathbb{N}}}~m}}
      \end{equation}
    \end{minipage}
    \hfill
    \strut
  \end{defn}
  which is also an equality relation.
\end{frame}

\begin{frame}
  \frametitle{Inductive / Coinductive definitions}
  \framesubtitle{Coinductive definition}
  We can define streams coinductively from the two constructors
  \begin{defn}[Stream]\setlengths\startwitheq
    \strut
    \hfill
    \begin{minipage}{0.275\linewidth}
      \begin{equation}
        \vcenter{\infer{\destructor{hd}~\term{s} : \type{A}}{\term{s} : \typeformer{stream}~\type{A}}}
      \end{equation}
    \end{minipage}
    \hfill
    \begin{minipage}{0.3\linewidth}
      \begin{equation}
        \vcenter{\infer{\destructor{tl}~\term{s}: \typeformer{stream}~\type{A}}{\term{s} : \typeformer{stream}~\type{A}}}
      \end{equation}
    \end{minipage}
    \hfill
    \strut
  \end{defn}
  for which we can coinductively define an equivalence relation
  \begin{defn}[Equivalence for Streams]\setlengths %\startwitheq
    \begin{equation}
      \vcenter{\infer{\term{s}~\relation{\sim_{stream}}~\term{t}}{\destructor{hd}~\term{s} = \destructor{hd}~\term{t} & \destructor{tl}~\term{s}~\relation{\sim_{stream}}~\destructor{tl}~\term{t}}}
    \end{equation}
  \end{defn}
  however, this bisimulation does not give us equality.
\end{frame}

\subsubsection{Univalent Foundations}
\begin{frame}
  \frametitle{Univalent Foundations}
  \framesubtitle{The Univalence Axiom}
  The univalence axiom says that equivalence is equivalent to equality
  \begin{defn}[Univalence]\setlengths\startwitheq
    \begin{equation}
      (\type{A} = \type{B}) \simeq (\type{A} \simeq \type{B})
    \end{equation}
  \end{defn}
  so if we work in a type theory where this holds, then \(\relation{\sim_{stream}}\) becomes an equality relation for streams.
\end{frame}

\subsubsection{Homotopy Type Theory (HoTT)}
\begin{frame}
  \frametitle{Homotopy Type Theory (HoTT)}
  % \framesubtitle{HoTT}
  Homotopy Type Theory (HoTT)
  \begin{itemize}
  \item is an intensional dependent type theory (builts on MLTT)
  \item assumes the univalence axiom
  \item has higher inductive types (HITs)
  \end{itemize}
  Types can be seen as spaces and elements as points of a space, meaning we get non-trivial equalities, that is equalities that are not reflexivity. As such we can construct higher inductive types, where we have both point and equality constructors.
\end{frame}

\begin{frame}[fragile]
  \frametitle{Cubical Type Theory (CTT)}
  HoTT is constructive, however the univalence axiom is not, since it is an axiom. We therefore use a Cubical Type Theory, where we can prove the univalence axiom and get constructivity. Cubical Type Theory is named after the composition principle, where given all but one side of a square, we get the entire square.
  \begin{figure}[h]
    \centering
    \begin{tikzcd}
      \type{A} \ar[r,"\pathterm{p} \ct \pathterm{q} \ct \pathterm{r}",dashed] & \type{B} \\
      \type{C} \ar[u,"\pathterm{p}^{-1}"] \ar[r,"\pathterm{q}"'] & \type{D} \ar[u,"\pathterm{r}"']
    \end{tikzcd}
  \end{figure}
  This work has been formalized using the Cubical Agda proof assistant.
\end{frame}

\begin{frame}
  \frametitle{H-Levels}
  HoTT is proof relevant, so we can have two proofs of a statement that might not be equal. Types in HoTT is classified in H-levels starting at -2.
  \begin{itemize}
  \item[\((-2)\)] Contractible types, with an element equal to all other elements
  \item[\((-1)\)] Mere propositions / hProp, all elements in a type are equal, but the type might not be inhabited.
  \item[\((0)\)] hSets, given two elements, then all equalities between the two elements are equal.
    \\
    \strut\hspace{-5mm}\(\vdots\)
  \item[\((n)\)] all equalities at H-level \((n+1)\) are equal.
  \end{itemize}
\end{frame}

\subsubsection{Propositional Truncation}
\begin{frame}
  \frametitle{Propositional Truncation}
  We can truncate types to H-level \((-1)\) with propositional truncation, defined as the following HIT
  \begin{defn}[Propositional Truncation]\setlengths\startwitheq
    \strut
    \hfill
    \begin{minipage}[b]{0.25\linewidth}
      \begin{equation}
        \vcenter{\infer{\constructor{|}\,\term{x}\,\constructor{|} : \| \type{A} \|}{\term{x} : \type{A}}}
      \end{equation}
    \end{minipage}
    \hfill
    \begin{minipage}[b]{0.4\linewidth}
      \begin{equation}
        \vcenter{\infer{\constructor{squash}~\term{x}~\term{y} : \term{x} \equiv \term{y}}{\term{x}, \term{y} : \|\type{A}\|}}
      \end{equation}
    \end{minipage}
    \hfill
    \strut
  \end{defn}
  which removes the proof relevance, since  when truncated all proofs of a statement only states that the type is inhabited.
\end{frame}

\subsubsection{Set Truncated Quotients}
\begin{frame}
  \frametitle{Set Truncated Quotients}
  We can define quotients as by the following HIT.
  \begin{defn}[Set truncated quotient]\setlengths \startwitheq
    \strut
    \hfill
    \begin{minipage}[b]{0.25\linewidth}
      \begin{equation}
        \vcenter{\infer{\quotientconstructor{\term{x}} : \type{A} / \relation{\mathcal{R}}}{\term{x} : \type{A}}}
      \end{equation}
    \end{minipage}
    \hfill
    \begin{minipage}[b]{0.4\linewidth}
      \begin{equation}
        \vcenter{\infer{\constructor{eq/}~\term{x}~\term{y}~\term{r} : \term{x} \equiv \term{y}}{\term{x}, \term{y} : \type{A} / \relation{\mathcal{R}} & \term{r} : \term{x}~\relation{\mathcal{R}}~\term{y}}}
      \end{equation}
    \end{minipage}
    \hfill
    \strut
    \begin{equation}
      \vcenter{\infer{\constructor{squash/} :\texttt{isSet}~(\type{A} / \relation{R})}{\strut}}
    \end{equation}
  \end{defn}
\end{frame}

\subsubsection{Axiom of Choice}
\begin{frame}
  \frametitle{Axiom of Choice}
  The axiom of choice in HoTT
  \begin{defn}[Axiom of Choice]\setlengths
    \begin{equation}
      \prod_{(\term{x} : \type{X})} \big\| \typeformer{Y}~\term{x} \big\| \rightarrow \bigg\| \prod_{(\term{x} : \type{X})} \typeformer{Y}~\term{x} \bigg\|
    \end{equation}
  \end{defn}
  and the axiom of countable choice
  \begin{defn}[Axiom of countable choice]\setlengths
    \begin{equation}
      \prod_{(n : \mathbb{N})} \big\| \typeformer{Y}~n \big\| \rightarrow \bigg\| \prod_{(n : \mathbb{N})} \typeformer{Y}~n \bigg\|
    \end{equation}
  \end{defn}
\end{frame}

\section{\texttt{M}-types}

\begin{frame}
  \frametitle{Overview}
  \tableofcontents[currentsection,subsubsectionstyle=hide]
\end{frame}

\subsection{Defining \texttt{M}-types}
\subsubsection{Containers and Polynomial functors}
\begin{frame}
  \frametitle{Containers and Polynomial functors}
  \begin{defn}\setlengths
    A Container (or signature) is a dependent pair \(\container{S} = \containerpair{\type{A}}{\typeformer{B}}\) for the types \(\type{A} : \universe{\mathcal{U}}\) and \(\typeformer{B} : \type{A} \rightarrow \universe{\mathcal{U}}\). 
  \end{defn}
  \begin{defn}\setlengths
    A polynomial functor \(\functor{P}_{\container{S}}\) (or extension) for a container \(\container{S} = \containerpair{\type{A}}{\typeformer{B}}\) is defined, for types as
    \begin{equation}
      \begin{aligned}
        &\functor{P}_{\container{S}} : \universe{\mathcal{U}} \rightarrow \universe{\mathcal{U}} \\
        &\functor{P}_{\container{S}}\,\type{X} = \sum_{(\term{a} : \type{A})} \typeformer{B}\term{a} \rightarrow \type{X}
      \end{aligned}
    \end{equation}
    and for a function \(\function{f} : \type{X} \rightarrow \type{Y}\) as
    \begin{equation}
      \begin{aligned}
        &\functor{P}_{\container{S}}\,\function{f} : \functor{P}_{\container{S}}\type{X} \rightarrow \functor{P}_{\container{S}}\type{Y} \\
        &\functor{P}_{\container{S}}\,\function{f}~(\term{a},\function{g}) = (\term{a}, \function{f} \circ \function{g}).
      \end{aligned}
    \end{equation}
  \end{defn}
\end{frame}

\subsubsection{Coalgebra}
\begin{frame}
  \frametitle{Coalgebra}
  \begin{defn}\setlengths
    A \(\functor{P}_{\container{S}}\)-coalgebra is defined as
    \begin{equation}
      \typeformer{Coalg}_{\container{S}} = \sum_{(\type{C} : \universe{\mathcal{U}})} \type{C} \rightarrow \functor{P}_{\container{S}}\type{C}.
    \end{equation}
    where we denote a \(\functor{P}_{\container{S}}\)-coalgebra given by a type \(\type{C}\) and function \(\function{\gamma}\) as \(\coalg{\type{C}}{\function{\gamma}}\). Coalgebras morphisms are defined as
    \begin{equation}
      \begin{aligned}
        \coalg{\type{C}}{\function{\gamma}} \Rightarrow \coalg{\type{D}}{\function{\delta}} &= \sum_{(\function{f} : \type{C} \rightarrow \type{D})} \function{\delta} \circ \function{f} = \functor{P} \function{f} \circ \function{\gamma}
      \end{aligned}
    \end{equation}
  \end{defn}
\end{frame}

\subsubsection{Finality and Requirements of M-types}
\begin{frame}
  \frametitle{Final Coalgebra}
  \begin{defn}[Final Coalgebra / \texttt{M}-type]\setlengths
    A final coalgebra \(\coalg{\type{X}}{\function{\rho}}\), is a coalgebra that fulfills
    \begin{equation}
      \typeformer{Final}_{\container{S}} := \sum_{(\coalg{\type{X}}{\function{\rho}} : \mathtt{Coalg}_{\container{S}})} \prod_{(\coalg{\type{C}}{\function{\gamma}} : \mathtt{Coalg}_{\container{S}})} \mathtt{isContr}~(\coalg{\type{C}}{\function{\gamma}} \Rightarrow \coalg{\type{X}}{\function{\rho}}).
    \end{equation}
    \label{defn:M-type-final-def}
    We define \texttt{M}-types as the coinductive type that fulfill the property
    \begin{equation}
      \prod_{(\coalg{\type{C}}{\function{\gamma}} : \mathtt{Coalg}_{\container{S}})} \mathtt{isContr}~(\coalg{\type{C}}{\function{\gamma}} \Rightarrow \coalg{\typeformer{M}_{\container{S}}}{\function{out}})
    \end{equation}
    We denote \texttt{M}-types as \(\typeformer{M}_{\containerpair{\type{A}}{\typeformer{B}}}\), \(\typeformer{M}_{\container{S}}\) or just \(\typeformer{M}\) when the container is clear from the context.
  \end{defn}
\end{frame}

\subsubsection{Chain and Limits}
\begin{frame}[fragile]
  \frametitle{Chain}
  \begin{defn}[Chain]\setlengths
    We define a chain as a family of morphisms \(\function{\pi}_{(n)} : \type{X}_{n+1} \rightarrow \type{X}_{n}\), over a family of types \(\type{X}_{n}\). See figure.
  \end{defn}
  \begin{figure}[h]
    \centering
    \begin{tikzcd}[row sep=large,column sep=large]
      \type{X}_0 & \type{X}_1 \ar[l,"\function{\pi}_{(0)}"] & \cdots \ar[l,"\function{\pi}_{(1)}"] & \type{X}_n \ar[l,"\function{\pi}_{(n-1)}"] & \type{X}_{n + 1} \ar[l,"\function{\pi}_{(n)}"] & \cdots \ar[l,"\function{\pi}_{(n+1)}"]
    \end{tikzcd}
  \end{figure}
  \begin{defn}\setlengths
    The limit of a chain is given as
    \begin{equation}
      \type{\mathcal{L}} = \sum_{(\term{x} : \prod_{(n : \mathbb{N})} \type{X}_n)} \prod_{(n : \mathbb{N})} (\function{\pi}_{(n)}~\term{x}_{n+1} \equiv \term{x}_n)
    \end{equation}
  \end{defn}
\end{frame}

\begin{frame}[fragile]
  \frametitle{Helper Lemmas}
  \begin{lem}[Limit Collapse]\setlengths
    \label{lem:limit-collapse}
    Given \(\function{\ell} : \prod_{(n : \mathbb{N})} (\type{X}_{n} \rightarrow \type{X}_{n+1})\) and \(\term{y} : \sum_{\left(\term{x} : \prod_{(n : \mathbb{N})} \type{X}_n\right)} \term{x}_{n + 1} \equiv \function{\ell}_n~\term{x}_n\) we get the equality \(\type{\mathcal{L}} \equiv \type{X}_0\).
  \end{lem}
  \begin{lem}[Cone and function to \texttt{M}-type equality]\setlengths
    \label{lem:function-to-M-type-is-cone}
    For all coalgebras \(\coalg{\type{C}}{\function{\gamma}}\) defined over the container \(\container{S}\), we get \(\type{C} \rightarrow \typeformer{M}_{\container{S}} \equiv \typeformer{Cone}_{\coalg{\type{C}}{\function{\gamma}}}\), where \(\typeformer{Cone} = \sum_{(\function{f} : \prod_{(n : \mathbb{N})} \type{C} \rightarrow \type{X}_n)} \prod_{(n : \mathbb{N})} \function{\pi}_{(n)} \circ \function{f}_{n+1} \equiv \function{f}_n\). (See figure)
  \end{lem}
  \begin{figure}[h]
    \centering
    \begin{tikzcd}[row sep=large,column sep=large]
      \type{X}_0 & \type{X}_1 \ar[l,"\function{\pi}_{(0)}"] & \cdots \ar[l,"\function{\pi}_{(1)}"] & \type{X}_n \ar[l,"\function{\pi}_{(n-1)}"] & \type{X}_{n + 1} \ar[l,"\function{\pi}_{(n)}"] & \cdots \ar[l,"\function{\pi}_{(n+1)}"] \\
      & & \type{C} \ar[ull,"\function{f}_0"] \ar[ul,"\function{f}_1"'] \ar[u,draw=none,"\cdots" description] \ar[ur,"\function{f}_n"] \ar[urr,"\function{f}_{n+1}"'] \ar[urrr,draw=none,"\cdots"'] & 
    \end{tikzcd}
  \end{figure}
\end{frame}

\begin{frame}<presentation:0>[noframenumbering] % proof frame hidden
  \frametitle{Helper Lemmas}
  \framesubtitle{Proof of Limit Collapse}
    \begin{proof}\setlengths
    We define the equality by the following isomorphism
    \begin{align}
      &\function{fun}_{\pathterm{\mathcal{L}collapse}}~(\term{x} , \pathterm{r}) = \term{x}_0 \\[-1mm]
      &\function{inv}_{\pathterm{\mathcal{L}collapse}}~\term{x}_0 = (\lambda\,n,~\function{\ell}^{(n)}~\term{x}_0)~,~(\lambda\,n,~\mathtt{refl}_{(\function{\ell}^{(n+1)}\,\term{x}_0)}) \\[-1mm]
      &\function{rinv}_{\pathterm{\mathcal{L}collapse}}~\term{x}_0 = \mathtt{refl}_{\term{x}_0}
    \end{align}
    where \(\function{\ell}^{(n)} = \function{\ell}_n \circ \function{\ell}_{n-1} \circ \dots \circ \function{\ell}_1 \circ \function{\ell}_0\). To define \(\function{linv}_{\pathterm{\mathcal{L}collapse}}~(\term{x} , \pathterm{r})\), we first define a fiber \((\type{X} , \term{z}, \function{\ell})\) over \(\mathbb{N}\) given some \(\term{z} : \type{X}_0\). Then any element of the type \(\sum_{\left(\term{x} : \prod_{(n : \mathbb{N})} \type{X}_n\right)} \term{x}_{n + 1} \equiv \function{\ell}_n~\term{x}_n\) is equal to a section over the fiber we defined. This means \(\term{y}\) is equal to a section. Since the sections are defined over \(\mathbb{N}\), which is an initial algebra for the functor \(\functor{G}\type{Y} = \unit + \type{Y}\), we get the sections are contractible. This means \(\term{y} \equiv \function{inv}_{\pathterm{\mathcal{L}collapse}} (\function{fun}_{\pathterm{\mathcal{L}collapse}}~\term{y})\) by uniqueness argument, since both sides of the equality are equal to sections over \(\mathbb{N}\).
  \end{proof}
\end{frame}

\begin{frame}<presentation:0>[noframenumbering] % proof frame hidden
  \frametitle{Helper Lemmas}
  \framesubtitle{Proof of Cone and function to \texttt{M}-type equality}
  \begin{proof}\setlengths
    We show the equality between \(\type{C} \rightarrow \typeformer{M}_{\container{S}}\) and \(\typeformer{Cone}_{\coalg{\type{C}}{\function{\gamma}}}\) as an isomorphism
    \begin{align}
      &\function{fun}_{\pathterm{collapse}}~\function{f} = (\lambda\,n\,\term{z},\pi_1~(\function{f}~\term{z})~n),(\lambda\,n~i~\term{a},~\pi_2~(\function{f}~\term{a})~n~i) \\
      &\function{inv}_{\pathterm{collapse}}~(\function{u},\function{q})~\term{z} = (\lambda\,n,~\function{u}~n~\term{z}),(\lambda\,n\,i,~\function{q}~n~i~\term{z}) \\
      &\function{rinv}_{\pathterm{collapse}}~(\function{u},\function{q}) = \mathtt{refl}_{(\function{u},\function{q})} \\
      &\function{linv}_{\pathterm{collapse}}~\function{f} = \mathtt{refl}_{\function{f}}
    \end{align}
  \end{proof}
\end{frame}

\subsubsection{Equality between \(\type{\mathcal{L}}\) and \(\functor{P}\,\type{\mathcal{L}}\)}
\begin{frame}
  \frametitle{Equality between \(\type{\mathcal{L}}\) and \(\functor{P}\,\type{\mathcal{L}}\)}
  \begin{thm}\setlengths
    Given a container \(\containerpair{\type{A}}{\typeformer{B}}\), we define a chain as the repeated application of \(\functor{P}\) to the unit element \(\type{X}_n = \functor{P}^n\,\unit\), and \(\function{\pi}_{(n)} = \functor{P}^n\,\constant{!}\) where \(\constant{!} : A \rightarrow \unit\) is the unique function into the unit type. Then there is an equality
    \begin{equation}
      \pathterm{shift} : \type{\mathcal{L}} \equiv \functor{P}\type{\mathcal{L}}
    \end{equation}
    where \(\type{\mathcal{L}}\) is the limit of this chain.
  \end{thm}
  \begin{block}{Proof structure}\setlengths
    The proof is done using the two helper lemmas
    \begin{align}
      \pathterm{\alpha} : \type{\mathcal{L}}^{\functor{P}} \equiv \functor{P}\,\type{\mathcal{L}} \\
      \pathterm{\mathcal{L}unique} : \type{\mathcal{L}} \equiv \type{\mathcal{L}}^{\functor{P}}
    \end{align}
    where \(\type{\mathcal{L}}^{\functor{P}}\) is the limit of the shifted chain defined as \(\type{X'}_n = \type{X}_{n+1}\) and \(\function{\pi'}_{(n)} = \function{\pi}_{(n+1)}\). With these two lemmas we get \(\pathterm{shift} = \pathterm{\alpha} \ct \pathterm{\mathcal{L}unique}\).
  \end{block}
\end{frame}

\begin{frame}
  \frametitle{Proof of uniqueness}
  \(\pathterm{\mathcal{L}unique}\) says the limit of a chain is unique.
  \begin{proof}\setlengths
  The proof of the equality \(\pathterm{\mathcal{L}unique}\) is given by the isomorphism
    \begin{align}
      \function{fun}_{\pathterm{\mathcal{L}unique}}~(\function{a} , \function{b}) &= \natcases{\unitelem}{\function{a}} , \natcases{\mathtt{refl}_{\unitelem}}{\function{b}} \\
    \function{inv}_{\pathterm{\mathcal{L}unique}}~(\function{a} , \function{b}) &= \function{a} \circ \constructor{succ}~,~\function{b} \circ \constructor{succ} \\
    \function{rinv}_{\pathterm{\mathcal{L}unique}}~(\function{a} , \function{b}) &= \mathtt{refl}_{(\function{a},\function{b})} \\
    \function{linv}_{\pathterm{\mathcal{L}unique}}~(\function{a} , \function{b}) &= \mathtt{refl}_{(\function{a},\function{b})}
    \end{align}
  \end{proof}
\end{frame}

\begin{frame}
  \frametitle{Proof of equality to shifted chain}
  \begin{proof}\setlengths
    The proof of \(\pathterm{\alpha}\) is given by the equalities
    \begin{align}
      \type{\mathcal{L}}^{\functor{P}} &\equiv \sum_{(\term{x} : \prod_{(n : \mathbb{N})} \type{X}_{n+1})} \prod_{(n : \mathbb{N})} \function{\pi}_{(n+1)}~\term{x}_{n+1} \equiv \term{x}_{n} \\
      &\equiv \hspace{-6mm} \sum_{(\term{x} : \prod_{(n : \mathbb{N})} \sum_{(\term{a} : \type{A})} \typeformer{B}\,\term{a} \rightarrow \type{X}_n)} \prod_{(n : \mathbb{N})} \function{\pi}_{(n+1)}~\term{x}_{n+1} \equiv \term{x}_{n} \label{eq:alpha-unfold}\\
      &\equiv \hspace{-10mm} \sum_{(\sum_{(\term{a} : \prod_{(n : \mathbb{N})} \type{A})} \prod_{(n : \mathbb{N})} \term{a}_{n+1} \equiv \term{a}_n)} \sum_{(\function{u} : \prod_{(n : \mathbb{N})} \typeformer{B}\,\term{a}_n \rightarrow \type{X}_n)} \prod_{(n : \mathbb{N})} \function{\pi}_{(n)} \circ \function{u}_{n+1} \equiv_* \function{u}_n \label{eq:alpha-swap}\\
      &\equiv \sum_{(\term{a} : \type{A})} \sum_{(\function{u} : \prod_{(n : \mathbb{N})} \typeformer{B}\,a \rightarrow \type{X}_n)} \prod_{(n : \mathbb{N})} \function{\pi}_{(n)} \circ \function{u}_{n+1} \equiv \function{u}_n \label{eq:alpha-limit-collapse} \\
      &\equiv \sum_{(\term{a} : \type{A})} \typeformer{B}\,\term{a} \rightarrow \type{\mathcal{L}} \label{eq:alpha-cone} \qquad \equiv \functor{P}\type{\mathcal{L}}
    \end{align}
    ~\\[-6mm]
  \end{proof}
  % TODO: Hints!
  % The equality from \eqref{eq:alpha-unfold} to \eqref{eq:alpha-swap} is rearranging terms, and adding a trivial point to the start of for the shifted chain. The equality from \eqref{eq:alpha-swap} and \eqref{eq:alpha-limit-collapse} is done with \lemref{limit-collapse} to collapse \(\sum_{(\term{a} : \mathbb{N} \rightarrow \type{A})} \prod_{(n : \mathbb{N})} \term{a}_{n+1} \equiv \term{a}_n\) to \(\type{A}\). For the equality between \eqref{eq:alpha-limit-collapse} and \eqref{eq:alpha-cone}, we use \lemref{function-to-M-type-is-cone}. The rest of the equalities are trivial.
\end{frame}

\begin{frame}
  \frametitle{In and Out of \texttt{M}-types}
  We define functions to construct and destruct \texttt{M}-types
  \begin{defn}\setlengths
    For the equality \(\pathterm{shift}\) we denote the functions back and forth as
    \begin{align}
      \function{out} &= \mathtt{transport}~\pathterm{shift} \\
      \function{in} &= \mathtt{transport}~\sym{\pathterm{shift}}.
    \end{align}
  \end{defn}
  This gives us the following diagram
  \begin{figure}[h]
    \centering
    \begin{tikzcd}
      \functor{P}\,\typeformer{M}_{\container{S}} \ar[d,shift left,"\function{in}"] \\
      \typeformer{M}_{\container{S}} \ar[u,shift left,"\function{out}"]
    \end{tikzcd}
  \end{figure}
\end{frame}

\subsubsection{Example: Streams}
\begin{frame}[fragile]
  \frametitle{Example: Streams}
  We can now define streams for a given type \(\type{A}\)
  \begin{defn}\setlengths
    We start with a container
    \begin{equation}
      \containerpair{\type{A}}{\unit}
    \end{equation}
    For which we get the polynomial functor
    \begin{equation}
      \functor{P}\,\type{X} = \type{A} \times \type{X}
    \end{equation}
    For which the we get the \texttt{M}-type for stream
    \begin{equation}
      \typeformer{Stream}~\type{A} = \type{A} \times \typeformer{Stream}~\type{A}
    \end{equation}
  \end{defn}
    \begin{figure}[h]
    \centering
    \begin{tikzcd}
      \type{A} \ar[dr,"\destructor{hd}"'] & \type{A} \times \typeformer{Stream}~\type{A} \ar[d,shift left,"\function{in}"] \ar[l,"\pi_1"'] \ar[r,"\pi_2"] & \typeformer{Stream}~\type{A} \ar[dl,"\destructor{tl}"] \\
      & \typeformer{Stream}~\type{A} \ar[u,shift left,"\function{out}"]
    \end{tikzcd}
  \end{figure}
\end{frame}

\subsubsection{\texttt{M}-types are final}
\begin{frame}
  \frametitle{\texttt{M}-types are final}
  \begin{thm}\setlengths
    The \texttt{M}-type \(\typeformer{M}_{\container{S}}\) is defined as the limit for a polynomial functor \(\functor{P}_{\container{S}}\). This definition fulfills the requirement that \(\typeformer{Final}_{\container{S}}~\type{\mathcal{L}}\).
  \end{thm}
  \begin{block}{Proof.}\setlengths
    By unfolding the definition, we need to show
    \begin{equation}
      \prod_{(\coalg{\type{C}}{\function{\gamma}} : \mathtt{Coalg}_{\container{S}})} \mathtt{isContr}~(\coalg{\type{C}}{\function{\gamma}} \Rightarrow \coalg{\type{\mathcal{L}}}{\function{out}})
    \end{equation}
    We do this by showing \((\coalg{\type{C}}{\function{\gamma}} \Rightarrow \coalg{\type{\mathcal{L}}}{\function{out}}) \equiv \unit\). We get
    \begin{align}
      &\equiv \sum_{(\function{f} : \type{C} \rightarrow \type{\mathcal{L}})} (\function{out} \circ \function{f} \equiv \functor{P}\function{f} \circ \function{\gamma}) \\[-2mm]
      &\equiv \sum_{(\function{f} : \type{C} \rightarrow \type{\mathcal{L}})} (\function{in} \circ \function{out} \circ \function{f} \equiv \function{in} \circ \functor{P}\function{f} \circ \function{\gamma}) \\[-2mm]
      &\equiv \sum_{(\function{f} : \type{C} \rightarrow \type{\mathcal{L}})} (\function{f} \equiv \function{in} \circ \functor{P}\function{f} \circ \function{\gamma})
    \end{align}
  \end{block}
\end{frame}

\begin{frame}[fragile]
  \begin{block}{Proof. (cont.)}\setlengths
    We now let \(\function{\psi} = \function{in} \circ \functor{P}\function{f} \circ \function{\gamma}\), which simplifies the expression to
    \begin{equation}
      \sum_{(\function{f} : \type{C} \rightarrow \type{\mathcal{L}})} (\function{f} \equiv \function{\psi}~\function{f})
    \end{equation}
    We then define \(\function{e}\) to be the function from right to left for the equality \(\type{C} \rightarrow \type{\mathcal{L}} \equiv \typeformer{Cone}_{\coalg{\type{C}}{\function{\gamma}}}\), this gives us the equality
    \begin{equation}
      \sum_{(\function{f} : \type{C} \rightarrow \type{\mathcal{L}})} (\function{f} \equiv \function{\psi}~\function{f}) \quad \equiv \quad \sum_{(\term{c} : \typeformer{Cone}_{\coalg{\type{C}}{\function{\gamma}}})} (\function{e}~\term{c} \equiv \function{\psi}~(\function{e}~\term{c}))
    \end{equation}
    By defining a function \(\function{\phi}~(\function{u},\function{g}) = (\function{\phi}_0~\function{u},\function{\phi}_1~\function{u}~\function{g})\) where
    \begin{align}
      \function{\phi}_0~\function{u} &= \natcases{(\lambda\,\_, \unitelem)}{\functor{P}\function{f} \circ \function{\gamma} \circ \function{u}} \\
      \function{\phi}_1~\function{u}~\function{g} &= \natcases{\mathtt{funExt}~(\lambda\,\_, \mathtt{refl}_{\unitelem})}{\mathtt{ap}~(\functor{P}\function{f} \circ \function{\gamma}) \circ \function{g}}
    \end{align}
    we get the commuting square\vspace{-4.5mm}
    \begin{figure}
      \centering
      \begin{tikzcd}
        \typeformer{Cone}_{\coalg{\type{C}}{\function{\gamma}}} \ar[d,"\function{\phi}"] \ar[r,"\function{e}"] & (\type{C} \rightarrow \type{\mathcal{L}}) \ar[d,"\function{\psi}"] \\
        \typeformer{Cone}_{\coalg{\type{C}}{\function{\gamma}}} \ar[r,"\function{e}"] & (\type{C} \rightarrow \type{\mathcal{L}})
      \end{tikzcd}
    \end{figure}
  \end{block}
\end{frame}

\begin{frame}
  \begin{block}{Proof. (cont.)}\setlengths
    We can then simplify to
    \begin{equation}
      \sum_{(\term{c} : \typeformer{Cone}_{\coalg{\type{C}}{\function{\gamma}}})} (\function{e}~\term{c} \equiv \function{e}~(\function{\phi}~\term{c}))
    \end{equation}
    We unfolding the definition of \(\function{\phi}\), to get the equalities
    \begin{align}
      &\quad~\sum_{(\term{c} : \typeformer{Cone}_{\coalg{\type{C}}{\function{\gamma}}})} (\term{c} \equiv \function{\phi}~\term{c}) \\
      &\equiv \sum_{(\term{(\function{u},\function{g})} : \typeformer{Cone}_{\coalg{\type{C}}{\function{\gamma}}})} ((\function{u},\function{g}) \equiv (\function{\phi}_0~\function{u},\function{\phi}_1~\function{u}~\function{g})) \\
      &\equiv \sum_{(\term{(\function{u},\function{g})} : \typeformer{Cone}_{\coalg{\type{C}}{\function{\gamma}}})} \sum_{(\pathterm{p} : \function{u} \equiv \function{\phi}_0\,\function{u})} \function{g} \equiv_{*} \function{\phi}_1~\function{u}~\function{g}
    \end{align}
    We rearrange the terms and unfold the definition of \(\typeformer{Cone}\) to get
    \begin{equation}
      \sum_{\big(\term{(\function{u},\pathterm{p})} : \sum_{(\function{u} : \prod_{(n:\mathbb{N})} \type{C} \rightarrow \type{X}_n)} \function{u} \equiv \function{\phi}_0\,\function{u}\big)} \sum_{(\function{g} : \prod_{(n:\mathbb{N})} \function{\pi}_{(n)} \circ \function{u}_{n+1} \equiv \function{u}_n)}\function{g} \equiv_{*} \function{\phi}_1~\function{u}~\function{g}
    \end{equation}
  \end{block}
\end{frame}

\begin{frame}
  \begin{block}{Proof. (cont.)}
  We define the following two equalities as instances of the limit collapse equality \(\type{\mathcal{L}} \equiv \type{X}_0\).
    \begin{align}
      \unit &\equiv \left(\sum_{(\function{u} : \prod_{(n:\mathbb{N})} \type{C} \rightarrow \type{X}_n)} \function{u} \equiv \function{\phi}_0\,\function{u}\right)  \\
      \unit &\equiv_{*} \left(\sum_{(\function{g} : \prod_{(n:\mathbb{N})} \function{\pi}_{(n)} \circ \function{u}_{n+1} \equiv \function{u}_n)}\function{g} \equiv_{*} \function{\phi}_1~\function{u}~\function{g}\right)
    \end{align}
    Using these two equalities, the proof simplifies to
    \begin{equation}
      \sum_{(\unitelem : \unit)} \unit \equiv \unit
    \end{equation}
    which is trivial. \qed
  \end{block}
\end{frame}

\subsection{Coinduction Principle}
\begin{frame}[fragile]
  \frametitle{Bisimulation}
  \begin{defn}\setlengths
    A relation \(\relation{\mathcal{R}} : \type{C} \rightarrow \type{C} \rightarrow \universe{\mathcal{U}}\) for a coalgebra \(\coalg{\type{C}}{\function{\gamma}} : \typeformer{Coalg}_{\container{S}}\), is a (strong) bisimulation relation if the type \(\type{\overline{\mathcal{R}}} = \sum_{(\term{a} : \type{C})} \sum_{(\term{b} : \type{C})} \term{a}~\relation{\mathcal{R}}~\term{b}\) and the function \(\function{\alpha}_{\relation{\mathcal{R}}} : \type{\overline{\mathcal{R}}} \rightarrow \functor{P}_{\container{S}}\,\type{\overline{\mathcal{R}}}\) forms a \(\functor{P}_{\container{S}}\)-coalgebra \(\coalg{\type{\overline{\mathcal{R}}}}{\function{\alpha}_{\relation{\mathcal{R}}}} : \mathtt{Coalg}_{\container{S}}\), making the diagram bellow commute (\(\Longrightarrow\) represents \(\functor{P}_{\container{S}}\)-coalgebra morphisms). That is
    \begin{equation}
      \function{\gamma} \circ \pi_1^{\type{\overline{\mathcal{R}}}} \equiv \functor{P}_{\container{S}}\,\pi_1^{\type{\overline{\mathcal{R}}}} \circ \function{\alpha}_{\relation{R}}
    \end{equation}
    and similarly for \(\pi_2^{\type{\overline{\mathcal{R}}}}\).
  \end{defn}
  \begin{figure}[h]
    \centering
    \begin{tikzcd}[row sep=huge,column sep=huge]
      \coalg{\type{C}}{\function{\gamma}} & \coalg{\type{\overline{\mathcal{R}}}}{\function{\alpha}_{\relation{\mathcal{R}}}} \ar[l,Rightarrow,"{\pi_1}^{\type{\overline{\mathcal{R}}}}"'] \ar[r,Rightarrow,"{\pi_2}^{\type{\overline{\mathcal{R}}}}"] & \coalg{\type{C}}{\function{\gamma}}
    \end{tikzcd}
  \end{figure}
\end{frame}

\begin{frame}[fragile]
  \frametitle{Coinduction Principle}
  \begin{thm}[Coinduction principle]\setlengths
    Given a relation \(\relation{\mathcal{R}}\), that is a bisimulation for a \texttt{M}-type, then (strongly) bisimilar elements \(\term{x}~\relation{\mathcal{R}}~\term{y}\) are equal \(\term{x} \equiv \term{y}\).
  \end{thm}
  \begin{proof}\setlengths
    Given a relation \(\relation{\mathcal{R}}\) that is bisimulation relation over a final \(\functor{P}\)-coalgebra \(\coalg{\typeformer{M}}{\function{out}}:\mathtt{Coalg}_{\container{S}}\) we get the diagram
    \begin{figure}[h]
      \centering
      \begin{tikzcd}[row sep=huge,column sep=huge]
        \coalg{\typeformer{M}}{\function{out}} & \coalg{\type{\overline{\mathcal{R}}}}{\function{\alpha}_{\relation{\mathcal{R}}}} \ar[l,Rightarrow,"{\pi_1}^{\type{\overline{\mathcal{R}}}}"'] \ar[r,Rightarrow,"{\pi_2}^{\type{\overline{\mathcal{R}}}}"] & \coalg{\typeformer{M}}{\function{out}}
      \end{tikzcd}
    \end{figure}
    By the finality of \(\coalg{\typeformer{M}}{\function{out}}\), we get a function \(\constant{!}\) from \(\typeformer{M}\) to \(\type{\overline{R}}\), which is unique, meaning \(\pi_1^{\type{\overline{\mathcal{R}}}} \equiv ~\constant{!} \equiv \pi_2^{\type{\overline{\mathcal{R}}}}\). Now given \(\term{r} : \term{x}~\relation{\mathcal{R}}~\term{y}\), we can construct the equality
    \begin{equation}
      \term{x} \equiv \pi_1^{\type{\overline{\mathcal{R}}}} (\term{x}, \term{y}, \term{r}) \equiv \pi_2^{\type{\overline{\mathcal{R}}}} (\term{x}, \term{y}, \term{r}) \equiv \term{y},
    \end{equation}
    giving us the coinduction principle for \texttt{M}-types.
  \end{proof}  
\end{frame}

\subsection{Examples of \texttt{M}-types}
\begin{frame}[fragile]
  \frametitle{Example: Delay Monad}
  We define a container
  \begin{equation}
    \containerpair{\type{R} + \unit}{[\empt, \unit]}
  \end{equation}
  and a polynomial functor
  \begin{equation}
    \functor{P}\,\type{X} = \sum_{(\term{x} : \type{R} + \unit)} \begin{cases} \empt & \term{x} = \mathtt{inl}~\term{r} \\ \unit & \term{x} = \mathtt{inr}~\unitelem \end{cases} \rightarrow \type{X},
  \end{equation}
  which simplifies to
  \begin{equation}
    \functor{P}_{\container{S}}\,\type{X} = \type{R} + \type{X}
  \end{equation}
  such that we get the diagram
  \begin{figure}[h]
    \centering
    \begin{tikzcd}[row sep=large,column sep=large]
      \type{R} \ar[dr,"\constructor{now}"'] \ar[r,"\texttt{inl}"] & \type{R} + \typeformer{Delay}~\type{R} \ar[d,shift left,"\function{in}"] & \typeformer{Delay}~\type{R} \ar[l,"\texttt{inr}"'] \ar[dl,"\constructor{later}"] \\
      & \typeformer{Delay}~\type{R} \ar[u,shift left,"\function{out}"]
    \end{tikzcd}
  \end{figure}
\end{frame}

\begin{frame}[fragile]
  \frametitle{Example: Event Trees}
  \begin{defn}\setlengths
    We define a container
    \begin{equation}
      \containerpair{\type{R} + \sum_{(\type{A} : \universe{\mathcal{U}})} \typeformer{E}\,\type{A}}{[\empt, \type{A}]}.
    \end{equation}
    and a polynomial functor which simplifies to
    \begin{equation}
      \functor{P}\,\type{X} = \type{R} + \sum_{(\type{A} : \universe{\mathcal{U}})} \typeformer{E}\,\type{A} \times (\type{A} \rightarrow \type{X}).
    \end{equation}
  \end{defn}
  \begin{figure}[h]
    \centering
    \begin{tikzcd}[row sep=large,column sep=small]
      \type{R} \ar[dr,"\constructor{leaf}"'] \ar[r,"\texttt{inl}"] & \type{R} + \sum_{(\type{A} : \universe{U})} \typeformer{E}\,\type{A} \times (\type{A} \rightarrow \typeformer{T}~\typeformer{E}~\type{R}) \ar[d,shift left,"\function{in}"]  & \sum_{(\type{A} : \universe{U})} \typeformer{E}\,\type{A} \times (\type{A} \rightarrow \typeformer{T}~\typeformer{E}~\type{R}) \ar[l,"\texttt{inr}"'] \ar[dl,"\constructor{node}"] \\
      & \typeformer{T}~\typeformer{E}~\type{R} \ar[u,shift left,"\function{out}"]
    \end{tikzcd}
  \end{figure}
\end{frame}

\begin{frame}
  \frametitle{Example: Interaction Trees (ITrees)}
  \begin{equation}
    \vcenter{\infer[\mathtt{I}_{\constructor{Ret}}]{\constructor{Ret}~\term{r} : \typeformer{itree}~\typeformer{E}\,\type{R}}{r : \type{R}}}
  \end{equation}
  \begin{equation}
    \vcenter{\infer[\mathtt{I}_{\constructor{Vis}}]{\constructor{Vis}~\term{a}~\term{f} : \typeformer{itree}~\typeformer{E}\,\type{R}}{\type{A} : \universe{\mathcal{U}} & \term{a} : \typeformer{E}\,\type{A} & \term{f} : \type{A} \rightarrow \typeformer{itree}~\typeformer{E}\,\type{R}}}.
  \end{equation}
  \begin{equation}
    \vcenter{\infer[\mathtt{E}_{\constructor{Tau}}]{\constructor{Tau}~\term{t} : \typeformer{itree}~\typeformer{E}\,\type{R}}{\term{t} : \typeformer{itree}~\typeformer{E}\,\type{R}}}.
  \end{equation}
  \begin{defn}\setlengths
    We define a container
    \begin{equation}
      \containerpair{\type{R}+\unit+\sum_{(\type{A}:\universe{\mathcal{U}})} (\typeformer{E}\,\type{A})}{[\empt , \unit , \type{A}]}.
    \end{equation}
    and a polynomial functor which simplifies to
    \begin{equation}
      \functor{P}\,\type{X} = \type{R} + \type{X} + \sum_{(\type{A} : \universe{\mathcal{U}})} (\typeformer{E}\,\type{A} \times (\type{A} \rightarrow \type{X}))
    \end{equation}
  \end{defn}
\end{frame}

\begin{frame}[shrink=36,fragile]
  \frametitle{Example: Interaction Trees (ITrees)}
  \vspace{35mm}
  \begin{figure}[h]
    \centering
    \begin{tikzcd}[row sep=large,column sep=large]
      \type{R} \ar[ddrr,"\constructor{Ret}"'] \ar[r,"\texttt{inl}\,\circ\,\texttt{inr}"] & \typeformer{IT}~\typeformer{E}~\type{R} + \type{R} + \sum_{(\type{A} : \universe{\mathcal{U}})} (\typeformer{E}\,\type{A} \times (\type{A} \rightarrow \typeformer{IT}~\typeformer{E}~\type{R})) \ar[ddr,"\function{in}"] & & \sum_{(\type{A} : \universe{U})} \typeformer{E}\,\type{A} \times (\type{A} \rightarrow \typeformer{IT}~\typeformer{E}~\type{R}) \ar[ll,"\texttt{inr}"'] \ar[ddl,"\constructor{Vis}"] \\
      & & \typeformer{IT}~\typeformer{E}~\type{R} \ar[ul,"\texttt{inl}\,\circ\,\texttt{inl}"'] \ar[d,"\constructor{Tau}"'] \\
      & & \typeformer{IT}~\typeformer{E}~\type{R}
    \end{tikzcd}
  \end{figure}
\end{frame}

\subsection{Rules for Construction M-types}
\begin{frame}
  \frametitle{Rules for Constructing \texttt{M}-types}
  Adding containers \(\containerpair{\type{A}}{\typeformer{B}}\) and \(\containerpair{\type{C}}{\typeformer{D}}\) for two constructors together is done by
\begin{equation}
  \containerpair{\type{A} + \type{C}}{ \begin{cases} \typeformer{B}\,\term{a} & \mathtt{inl}~\term{a} \\ \typeformer{D}\,\term{c} & \mathtt{inr}~\term{c} \end{cases}}
\end{equation}
whereas adding containers for two destructors is done by
\begin{equation}
  \containerpairsimple{\type{A} \times \type{C}}{\lambda\,(\term{a},\term{c}),~\typeformer{B}\,\term{a} + \typeformer{D}\,\term{c}}  
\end{equation}
However combining both destructors and constructors is not as simple, which is also the case if we use records or other constructions to define such types.
\end{frame}

\begin{frame}
  \frametitle{Rules for Constructing \texttt{M}-types}
  Given a record, which is a list of fields \(\term{f}_1 : \type{F}_1,\term{f}_2 : \type{F}_2,\dots,\term{f}_n : \type{F}_n\), we can construct the \texttt{M}-types by the container
\begin{equation}
  \containerpairsimple{\type{F}_1 \times \type{F}_2 \times \dots \times \type{F}_n}{\empt}
\end{equation}
where each destructor \(\destructor{d_n} : \type{T} \rightarrow \type{F_n}\) for the field \(\term{f}_n\) will be defined as \(\destructor{d_n}~\term{t} = \pi_n~(\function{out}~\term{t})\). However fields in a coinductive record may depend on previous defined fields.
\begin{equation}
  \containerpairsimple{\sum_{(\term{f}_1 : \type{F}_1)} \sum_{(\term{f_2} : \typeformer{F}_2)} \cdots \sum_{(\term{f}_{n-1} : \typeformer{F}_{n-1})}\typeformer{F}_n}{\empt}
\end{equation}
even in a dependent way
\begin{equation}
  \containerpair{\sum_{(\type{f}_1 : \universe{\mathcal{U}})} \left(\unit \times \sum_{(\term{f}_3 : \typeformer{F}_3)} \cdots \sum_{(\term{f}_{n-1} : \typeformer{F}_{n-1})}\typeformer{F}_n\right)}{\lambda\,(\type{f}_1,\unitelem,\term{f}_3,\dots),\typeformer{F}_2}
\end{equation}
\end{frame}

\begin{frame}
  \frametitle{Rules for Constructing \texttt{M}-types}
  Another case is that the type of a field is dependent, but by currying we get
  \begin{equation}
    \containerpair{\sum_{(\type{f}_1 : \universe{\mathcal{U}})} \sum_{(\term{f}_3 : \typeformer{F_3})} \cdots \sum_{(\term{f}_{n-1} : \typeformer{F_{n-1}})}\typeformer{F_n}}{\lambda\,(\type{f}_1,\term{f}_3,\dots),\sum_{(\term{x} : \type{f}_1)} (\typeformer{B}\,\term{x})}
  \end{equation}
  so we expect that a type defined as a (coinductive) record is equal to the version defined as a \texttt{M}-type. However, this has not been shown formally.
\end{frame}

\section{Q\texttt{M}-types}

\begin{frame}
  \frametitle{Overview}
  \tableofcontents[currentsection,subsubsectionstyle=hide]
\end{frame}

\begin{frame}
  \frametitle{Q\texttt{M}-types}
  We will start by looking at two ways of quotienting \texttt{M}-types. \\
  Set truncated quotients and  quotient inductive-inductive type (QIIT). A QIIT is a type that is defined at the same time as a relation over that type and then set truncated.
\end{frame}

\begin{frame}
  \frametitle{Multiset}
  We construct the type of multisets as a Q\texttt{M}-type and as a QIIT and show these are equal. The base type is given by
  \begin{defn}\setlengths
    We define \(\type{A}\)-valued \(\mathbb{N}\)-branching trees \(\typeformer{T}\,\type{A}\) as the \texttt{M}-type defined by the container
    \begin{equation}
      \containerpair{\type{A}+\unit}{[\empt , \mathbb{N}]}
    \end{equation}
    For which we get the constructors\\
      \strut
      \hfill
      \begin{minipage}[b]{0.30\linewidth}
        \begin{equation}
          \vcenter{\infer{\constructor{leaf}~\term{a} : \typeformer{T}\,\type{A}}{\term{a} : \type{A}}}
        \end{equation}
      \end{minipage}
      \hfill
      \begin{minipage}[b]{0.30\linewidth}
        \begin{equation}
          \vcenter{\infer{\constructor{node}~\function{f} : \typeformer{T}\,\type{A}}{\function{f} : \mathbb{N} \rightarrow \typeformer{T}\,\type{A}}}
        \end{equation}
      \end{minipage}
      \hfill
      \strut
    \end{defn}
    The QIIT is defined by the point constructors \(\constructor{leaf}\) and \(\constructor{node}\), and equality constructors
    \begin{equation}
      \vcenter{\infer[\constructor{perm}]{\constructor{node}~\function{f}~\equiv~\constructor{node}~(\function{f} \circ \function{g})}{\function{f} : \mathbb{N} \rightarrow \typeformer{T}\,\type{A} & \function{g} : \mathbb{N} \rightarrow \mathbb{N} & \mathtt{isIso}~\function{g}}}
    \end{equation}
    and \(\constructor{MS\text{-}isSet}\) which set truncates the type.    
\end{frame}

\subsection{Multisets}
\begin{frame}
  \frametitle{Multisets}
  We get the Q\texttt{M}-type for multisets by quotienting the tree type by the relation defined by the constructors\\
  \strut
  \hfill
  \begin{minipage}{0.35\linewidth}
    \begin{equation}
      \vcenter{\infer[\constructor{\sim_{leaf}}]{\constructor{leaf}~\term{x}~\relation{\sim_{T}}~\constructor{leaf}~\term{y}}{\term{x} \equiv \term{y}}}
    \end{equation}
  \end{minipage}
  \hfill
  \begin{minipage}{0.35\linewidth}
    \begin{equation}
      \vcenter{\infer[\constructor{\sim_{node}}]{\constructor{node}~\function{f}~\relation{\sim_T}~\constructor{node}~\function{g}}{\prod_{(n : \mathbb{N})} \function{f}~n~\relation{\sim_T}~\function{g}~n}}
    \end{equation}
  \end{minipage}
  \hfill
  \strut
  \begin{equation}
    \vcenter{\infer[\constructor{\sim_{perm}}]{\constructor{node}~\function{f}~\relation{\sim_T}~\constructor{node}~(\function{f} \circ \function{g})}{\function{f} : \mathbb{N} \rightarrow \typeformer{T}\,\type{A} & \function{g} : \mathbb{N} \rightarrow \mathbb{N} & \mathtt{isIso}~\function{g}}}
  \end{equation}
\end{frame}

\begin{frame}
  \frametitle{Multisets}
  We define a function between the two definitions
    \begin{equation}
    \begin{aligned}
      &\function{T{\rightarrow}MS}~(\constructor{leaf}_{\typeformer{T}}~\term{x}) = \constructor{leaf}_{\typeformer{MS}}~\term{x} \\
      &\function{T{\rightarrow}MS}~(\constructor{node}_{\typeformer{T}}~\function{f}) = \constructor{node}_{\typeformer{MS}}~(\function{T{\rightarrow}MS} \circ \function{f})
    \end{aligned}
  \end{equation}
  \begin{lem}
    If \(\term{x}, \term{y} : \typeformer{T}\,\type{A}\) are weakly bisimilar \(\term{p} : \term{x}~\relation{\sim_{T}}~\term{y}\) then \(\function{T{\rightarrow}MS}~\term{x} \equiv \function{T{\rightarrow}MS}~\term{y}\).
  \end{lem}
  \begin{proof}\setlengths
    We do the proof by casing on the constructor for the weak bisimilarity.
    \begin{equation}
      \begin{aligned}
        &\function{T{\rightarrow}MS\text{-}{\sim}{\rightarrow}{\equiv}}~(\constructor{\sim_{leaf}}~\pathterm{p}) = \mathtt{ap}~\constructor{leaf}_{\typeformer{MS}}~\pathterm{p} \\
        &\function{T{\rightarrow}MS\text{-}{\sim}{\rightarrow}{\equiv}}~(\constructor{\sim_{node}}~\function{k}) = \mathtt{ap}~\constructor{node}_{\typeformer{MS}}~(\mathtt{funExt}~(\function{T{\rightarrow}MS\text{-}{\sim}{\rightarrow}{\equiv}} \circ \function{k})) \\
        &\function{T{\rightarrow}MS\text{-}{\sim}{\rightarrow}{\equiv}}~(\constructor{\sim_{perm}}~\function{f}~\function{g}~\term{e}) = \constructor{perm}~(\function{T{\rightarrow}MS} \circ \function{f})~\function{g}~\term{e}
      \end{aligned}
    \end{equation}
  \end{proof}
\end{frame}

\begin{frame}
  \frametitle{Multisets}
  With this lemma, we can lift the function \(\function{T{\rightarrow}MS}\) to the quotient, using the recursor for quotients. We now show that this function is injective and surjective. We start with a helper lemma.
  \\ \\
  \begin{lem}
    Given an equality \(\function{T{\rightarrow}MS}~\term{x} \equiv \function{T{\rightarrow}MS}~\term{y}\) then \(\term{x}~\relation{\sim_{T}}~\term{y}\).
  \end{lem}
  \begin{proof}
    If \(\term{x}\) and \(\term{y}\) are leaves with values \(\term{a}\) and \(\term{b}\) then \(\term{a} \equiv \term{b}\) by the injectivity of the constructor \(\constructor{leaf}_{\typeformer{MS}}\), giving us the wanted bisimilarity by \(\constructor{\sim_{leaf}}\). If \(\term{x}\) and \(\term{y}\) are nodes defined by functions \(\function{f}\) and \(\function{g}\), then by the injectivity of the constructor \(\constructor{node}_{\typeformer{MS}}\), we get \(\prod_{(n : \mathbb{N})} \function{f}~n \equiv \function{g}~n\) and by induction we get \(\prod_{(n : \mathbb{N})} \function{f}~n~\relation{\sim_{T}}~\function{g}~n\), making us able to use \(\constructor{\sim_{node}}\) to construct the bisimilarity. We do not get any other equalities, since \(\constructor{leaf}_{\typeformer{MS}}\) and \(\constructor{node}_{\typeformer{MS}}\) are disjoint.
  \end{proof}
\end{frame}

\begin{frame}
  \frametitle{Multiset}
  \begin{lem}
    The function \(\function{T/{\sim}{\rightarrow}MS}\) is injective, meaning \\
    \(\function{T/{\sim}{\rightarrow}MS}~\term{x} \equiv \function{T/{\sim}{\rightarrow}MS}~\term{y}\) implies \(\term{x} \equiv \term{y}\).
  \end{lem}\setlengths
  \begin{proof}
    We show injectivity by doing propositional elimination, with
    \begin{equation}
      \typeformer{P} = (\lambda\,\term{x},~\function{T/{\sim}{\rightarrow}MS}~\term{x} \equiv \function{T/{\sim}{\rightarrow}MS}~\term{y} \rightarrow \term{x} \equiv \term{y})
    \end{equation}
    then
    \begin{equation}
      \function{elimProp}~(\lambda\,\term{a},~\function{elimProp} (\lambda\,\term{b},~\constructor{eq/}~\term{a}~\term{b} \circ \function{T{\rightarrow}MS\text{-}{\equiv}{\rightarrow}{\sim}}~\term{a}~\term{b})~\term{y})~\term{x}
    \end{equation}
  \end{proof}
\end{frame}

\begin{frame}
  \frametitle{Multiset}
  \begin{defn}
    We define the propositional eliminator for \(\typeformer{MS}\).
    \begin{equation}
      \begin{aligned}
        &\function{elimProp_{MS}}~(\constructor{leaf}~\term{x}) = \function{P_{leaf}}~\term{x} \\
        &\function{elimProp_{MS}}~(\constructor{node}~\term{f}) = \function{P_{node}}~(\function{elimProp_{MS}} \circ f) \\[1mm]
        &\function{elimProp_{MS}}~(\constructor{perm}~\function{f}~\function{g}~\term{e}~i) = \\[-1mm]
        &\quad\function{isProp{\rightarrow}isDepProp}~\function{P_{prop}}~\\
        &\quad\qquad(\function{elimProp_{MS}}~(\constructor{node}~\function{f}))~(\function{elimProp_{MS}}~(\constructor{node}~(\function{f} \circ \function{g})))\\[-1mm]
        &\quad\qquad(\constructor{perm}~\function{f}~\function{g}~\term{e})~i \\[1mm]
        &\function{elimProp_{MS}}~(\constructor{MS\text{-}isSet}~\term{a}~\term{b}~\pathterm{p}~\pathterm{q}~i~j) = \\[-1mm]
        &\quad\function{isSet{\rightarrow}isDepSet}~(\function{isProp{\rightarrow}isSet} \circ \function{P_{prop}})~\\
        &\quad\qquad(\function{elimProp_{MS}}~\term{a})~(\function{elimProp_{MS}}~\term{b})\\[-1mm]
        &\quad\qquad(\mathtt{ap}~\function{elimProp_{MS}}~\pathterm{p})~(\mathtt{ap}~\function{elimProp_{MS}}~\pathterm{q})\\
        &\quad\qquad(\constructor{MS\text{-}isSet}~\term{a}~\term{b}~\pathterm{p}~\pathterm{q})~i~j \\
      \end{aligned}
    \end{equation}
  \end{defn}
\end{frame}

\begin{frame}
  \frametitle{Multiset}
  \begin{lem}
    The function \(\function{T/{\sim}{\rightarrow}MS}\) is surjective, assuming the axiom of countable choice. Being surjective means for all \(\term{b} : \typeformer{MS}\,\type{X}\) we have \(\| \Sigma_{(\term{x} : \typeformer{T}\,\type{X} / \relation{\sim_{T}})}~\function{T/{\sim}{\rightarrow}MS}~\term{x} \equiv \term{b} \|\).
  \end{lem}
    \begin{block}{Proof.}\setlengths
      We do propositional elimination of \(\typeformer{MS}~\type{X}\), with
      \begin{equation}
        \typeformer{P}~\term{y} = \| \Sigma_{(\term{x} : \typeformer{T}\,\type{X} / \relation{\sim_{T}})}~\function{T/{\sim}{\rightarrow}MS}~\term{x} \equiv \term{y} \|
      \end{equation}
      For the leaf case, we have the simple equality
      \begin{equation}
        \function{T/{\sim}{\rightarrow}MS}~\quotientconstructor{\constructor{leaf}_{\typeformer{T}}~\term{x} / \relation{\sim_{T}}} \equiv \constructor{leaf}_{\typeformer{MS}}~\term{x}
      \end{equation}
      For the node case with \(\constructor{node}~\function{f}\), we get the following function by induction
    \begin{equation}
      \function{f'} : \prod_{(n : \mathbb{N})} \bigg\| \sum_{(\term{y} : \typeformer{T}\,\type{X} / \relation{\sim_{T}})} \function{T/{\sim}{\rightarrow}MS}~\term{y} \equiv \function{f}\,n \bigg\|
    \end{equation}
  \end{block}
\end{frame}

\begin{frame}
  \frametitle{Multiset}
  \begin{block}{Proof. (cont.)}\setlengths
    Using the axiom of countable choice, we get
    \begin{equation}
      \bigg\| \prod_{(n : \mathbb{N})} \sum_{(\term{y} : \typeformer{T}\,\type{X} / \relation{\sim_{T}})} \function{T/{\sim}{\rightarrow}MS}~\term{y} \equiv \function{f}\,n \bigg\|
    \end{equation}
    By the surjectivity of \(\quotientconstructor{{\cdot}}\) we get
    \begin{equation}
      \bigg\| \prod_{(n : \mathbb{N})} \sum_{(\term{y} : \typeformer{T}\,\type{X})} \function{T{\rightarrow}MS}~\term{y} \equiv \function{f}\,n \bigg\|
    \end{equation}
    By swapping \(\Pi\) and \(\Sigma\) and functional extensionality we get
    \begin{equation}
      \bigg\| \sum_{(\term{g} : \mathbb{N} \rightarrow \typeformer{T}\,\type{X})} \function{T{\rightarrow}MS} \circ \function{g} \equiv \function{f} \bigg\|
    \end{equation}
    which gives us the equalities
    \begin{equation}
      \function{T/{\sim}{\rightarrow}MS}~\quotientconstructor{\constructor{node}~\function{g}} \equiv \constructor{node}~\function{f}
    \end{equation}
    completing the case for the node constructor and thereby the proof. \qed
  \end{block}
\end{frame}

\subsection{Partiality Monad}
\begin{frame}
  \frametitle{Partiality Monad}
  \textcolor{red}{TODO TODO TODO TODO TODO TODO TODO TODO TODO TODO TODO TODO TODO TODO TODO TODO TODO TODO TODO TODO TODO TODO TODO TODO TODO TODO TODO TODO TODO TODO TODO TODO TODO TODO TODO TODO TODO TODO TODO TODO TODO TODO TODO TODO TODO TODO TODO TODO TODO TODO TODO TODO TODO TODO TODO TODO TODO TODO TODO TODO TODO TODO TODO TODO TODO TODO TODO TODO TODO TODO TODO TODO TODO TODO TODO TODO TODO TODO TODO TODO TODO TODO TODO TODO TODO TODO TODO TODO TODO TODO TODO TODO TODO TODO TODO TODO TODO TODO TODO TODO TODO TODO TODO TODO TODO TODO TODO TODO TODO TODO TODO TODO TODO TODO TODO TODO TODO TODO TODO TODO TODO TODO TODO TODO TODO TODO TODO TODO TODO TODO TODO TODO TODO TODO TODO TODO TODO TODO TODO TODO TODO TODO TODO TODO}
\end{frame}

\subsection{Alternative Quotient \texttt{M}-type (QPF)}
\begin{frame}
  \frametitle{Quotient Polynomial Functor (QPF)}
  We can define quotiented \texttt{M}-types from a quotient polynomial functor.
  \begin{defn}[Quotient Polynomial Functor]\setlengths
    We define a quotient polynomial functor (QPF), for types as
    \begin{equation}
      \functor{F}\,\type{X} =  \sum_{(\term{a} : \type{A})} ((\typeformer{B}\,\term{a} \rightarrow \type{X}) / \relation{\sim}_{\term{a}})
    \end{equation}
    and for a function \(\function{f} : \type{X} \rightarrow \type{Y}\), we use the quotient eliminator with
    \begin{equation}
      \typeformer{P} = \lambda\,\_, (\typeformer{B}\,\term{a} \rightarrow \type{Y})/\relation{\sim}_{\term{a}}
    \end{equation}
    The base case is \(\function{f} = \lambda\,\_,\quotientconstructor{\function{f} \circ \function{g}}\) and equality case is \(\function{f_{eq}} = \lambda\,\function{x}\,\function{y}\,\term{r},~\constructor{eq/}~(\function{f} \circ \function{x})~(\function{f} \circ \function{y})~(\function{\sim_{ap}}~\function{f}~\term{r})\), where \(\function{\sim_{ap}}\) says that given a function \(\function{f}\) and \(\function{x}~\relation{\sim}_{\term{a}}~\function{y}\) then \(\function{f} \circ \function{x}~\relation{\sim}_{\term{a}}~\function{f} \circ \function{y}\). With this the definition for the quotient polynomial functor for functions is
    \begin{equation}
      \functor{F}\,\function{f}~(\term{a} , \function{g}) = (\term{a} , \function{elim}~\function{g})
    \end{equation}
    completing the definition of a quotiented polynomial functor.
  \end{defn}
\end{frame}

\begin{frame}[fragile]
  \frametitle{Alternative Construction of Q\texttt{M}-types}
  If there is a function \(\function{abs}_{\type{X}} : \functor{P}\,\type{X} \rightarrow \functor{F}\,\type{X}\) that makes the diagram bellow commute, then we can construct the final \(\functor{F}\)-coalgebra, to get another notion of quotiented \texttt{M}-type.
  \begin{figure}[h]
    \centering
    \begin{tikzcd}
      \functor{P}\,\type{X} \ar[r,"\functor{P}\,\function{f}"] \ar[d,"\function{abs}_{\type{X}}"'] & \functor{P}\,\type{Y} \ar[d,"\function{abs}_{\type{Y}}"] \\
      \functor{F}\,\type{X} \ar[r,"\functor{F}\,\function{f}"] & \functor{F}\,\type{Y}
    \end{tikzcd}
  \end{figure}
  % If \(\function{abs}\) is surjective, then the square will commute. % TODO: <- NOTE
  We get a surjective function \(\function{QM\text{-}intro} : \typeformer{M}_{\container{S}} \rightarrow \typeformer{M}_{\container{S}} / \relation{\sim}\) from this construction. The function is given as \(\function{abs}_{\typeformer{M}}\), which is the instantiation of \(\function{abs}\) between the limits of the two sequences. 
\end{frame}

\begin{frame}
  \frametitle{Construction QIITs}
  We can construct a QIIT from a M-type and a relation, by defining them simultanouosly, and we can then show, it equal to the Q\texttt{M}-type by giving
  \begin{itemize}
  \item A function \function{QM{\rightarrow}QIIT} from Q\texttt{M} to QIIT, taking the Q\texttt{M} constructors to their corresponding QIIT constructor.
  \item A proof \(\function{QM{\rightarrow}QIIT\text{-}{\sim}{\rightarrow}{\equiv}}\) that \(\term{x}~\relation{\sim_{QM}}~\term{y}\) implies \(\function{QM{\rightarrow}QIIT}~\term{x} \equiv \function{QM{\rightarrow}QIIT}~\term{y}\), which should follow by the constructors we added to the QIIT for the relation.
  \item Lifting \(\function{QM{\rightarrow}QIIT}\) to \(\function{QM/{\sim}{\rightarrow}QIIT}\) using the quotient recursor with \(\function{QM{\rightarrow}QIIT\text{-}{\sim}{\rightarrow}{\equiv}}\) and \(\constructor{QIIT\text{-}isSet}\)
  \item Showing injectivity by using propositional elimination of the quotient together with the inverse of \(\function{QM{\rightarrow}QIIT\text{-}{\sim}{\rightarrow}{\equiv}}\), namely \(\function{QM{\rightarrow}QIIT\text{-}injective}\) saying that \(\function{QM{\rightarrow}QIIT}~\term{x} \equiv \function{QM{\rightarrow}QIIT}~\term{y}\) implies \(\term{x}~\relation{\sim_{QM}}~\term{y}\), which again should follow from how the relations constructors were defined.
  \item Lastly we show surjectivity by induction using the eliminator of QIIT and the axiom of choice and the surjectivity of \(\quotientconstructor{{\cdot}}\).
  \end{itemize}  
\end{frame}

\begin{frame}
  However this is not enough...
\end{frame}

\section{Conclusion}
\begin{frame}
  \frametitle{Conclusion}
\end{frame}

\begin{frame}
  \frametitle{Future Work}
\end{frame}

\end{document}
